\documentclass[10pt,a4paper]{amsart}
\usepackage[margin=19mm]{geometry}
\usepackage{amsmath,amssymb,mathtools}
\usepackage[scr=rsfs,cal=euler]{mathalfa}
\usepackage{xfrac}
\usepackage[unicode,hidelinks,bookmarks=false]{hyperref}
\newcommand{\ie}{i.e.\ }
\newcommand{\cf}{cf.\ }
\newcommand{\resp}{respectively\ }
\newcommand{\Subsection}{\subsection}
\providecommand{\nobreakdash}{\nobreak\hbox{-}}
\setlength{\emergencystretch}{2em}
\multlinegap0pt
\usepackage{bm}
\usepackage{amssymb}
\usepackage{mathtools}
\usepackage[all]{xy}
\usepackage[shortlabels]{enumitem}
\newtheorem{dfn}{Definition}
\newtheorem{thm}{Theorem}
\newtheorem{thmlist}{Thmlist}
\newtheorem{lem}{Lemma}
\newtheorem{proposition}{Proposition}
\usepackage{cleveref}
\crefname{dfn}{Definition}{Definitions}
\crefname{thm}{Theorem}{Theorems}
\crefname{thmlist}{Thmlist}{Thmlists}
\crefname{lem}{Lemma}{Lemmas}
\crefname{proposition}{Proposition}{Propositions}
\newcommand{\GL}{{\rm GL}}
\newcommand{\bC}{\mathbb{C}}
\newcommand{\bF}{\mathbb{F}}
\newcommand{\bN}{\mathbb{N}}
\newcommand{\cL}{\mathcal L}
\newcommand{\sD}{\mathscr{D}}
\title{Topology, Hyperbolicity, and the Shafarevich Conjecture for Complex Algebraic Varieties}
\author{Ya Deng}
\date{Source: arXiv:2512.24458v1 (CC BY 4.0)}
\begin{document}
\maketitle

\noindent\emph{Source and licence.} Topology, Hyperbolicity, and the Shafarevich Conjecture for Complex Algebraic Varieties. Version arXiv:2512.24458v1 (CC BY 4.0), distributed by arXiv under the Creative Commons Attribution 4.0 International licence, http://creativecommons.org/licenses/by/4.0/, as recorded in the arXiv licence metadata for this exact version and shown on the abstract page https://arxiv.org/abs/2512.24458v1. Full licence notice: \texttt{licenses/arxiv-2512.24458-CC-BY-4.0.txt}.\par\smallskip

\noindent\textit{Editorial scope.} Counted unit: the article's thm thm:sym (source0.tex lines 1854--1861). The article's complete original proof of it is delivered with it (source0.tex lines 1863--1933, one \texttt{proof} environment of 71 lines, which the article places away from the statement). No mathematics is written, repaired or re-derived: the statement and the proof are the article's own lines, in the article's own notation, and no retained line is substituted or re-flowed.  Retained uncounted as the source-local context the proof needs: lines 1537--1555; lines 1580--1589; lines 1601--1603; lines 1672--1674; lines 1760--1762; lines 1787--1790.  Nothing was omitted from any retained span, so the package carries no omission record.  The retained lines cite 8 works, whose entries are transcribed verbatim from the article's own bibliography source in the e-print and delivered with short editorial labels recorded in the item's reference mapping.  No retained line contains a figure environment, an image-inclusion command or drawing code, so no image is delivered and the package declares no asset dependency.  Editorial formatting only, in addition: an AMS \texttt{amsart} preamble that restates the article's own declarations and macros used by the retained lines, namely the environments \texttt{dfn}, \texttt{thm}, \texttt{thmlist}, \texttt{lem}, \texttt{proposition} on separate counters, together with the standard packages those lines need.  The retained environments print in this document as Dfn 1, Theorem 1, Thmlist 1, Lemma 1, Proposition 1 in order of appearance, whereas the article prints the same items with the numbering of its own full document.  The complete native source of the article as distributed in the arXiv e-print for this exact version is bundled unchanged as \texttt{sources/191899/source0.tex}.
%
Furthermore, using the Lipschitz property of the harmonic map, one can show that $\eta_u$ extends to a multivalued 1-form on all of $X$. Since $u$ has logarithmic energy growth, in \cite{BDDM} we can prove that  $\eta_u$ extends to a multivalued section for the logarithmic cotangent bundle $\Omega_{\overline{X}}(\log D)$, where $\overline{X}$ is a smooth projective compactification for $X$ and $D:=\overline{X}\backslash X$ is a simple normal crossing divisor.     We refer the reader to \cite[\S 3.1]{CDY22} for the formal definition of \emph{multivalued sections} of a holomorphic vector bundle. Less formally, we have
\begin{dfn}[Multivalued  section]\label{def:multivalued}
	Let $X$ be a complex manifold and let $E$ be a  holomorphic  vector bundle on $X$.  A \emph{multivalued section} $\eta$ on $X$ is   a formal sum $ \sum_{i=1}^{m}n_i\Gamma_i$, where each $n_i\in \mathbb{Z}_{>0}$ and each $\Gamma_i$  is an irreducible closed subvariety of $E$ such that  
	the natural projection $\Gamma_i\to X$ is a finite surjective morphism.  
\end{dfn}
In \cite{DM24},  Mese and the author established the uniqueness of harmonic maps in a suitable setting. In particular, we showed that $\eta_u$ is independent of the choice of $u$. Therefore, we denote this form by $\eta_\varrho$.   It also extends to a multivalued section of the logarithmic bundle $\Omega_{\overline{X}}(\log D)\to \overline{X}$. 
\begin{dfn}[Multivalued 1-form]\label{def:multivaluedform}
	The multivalued 1-form $\eta_\varrho$ on $X$ described above is called the \emph{multivalued 1-form associated with $\varrho$}.
\end{dfn}
On the other hand, we can construct an analytic object associated with $\varrho$ as follows.   For the open neighborhood  $\Omega_x$ of $x$ as above, we define  a smooth, real, semi-positive $(1,1)$-form by
\begin{align}\label{eq:form}
	\sqrt{-1}\sum_{i=1}^{N}\partial u_{A,i} \wedge \bar{\partial} u_{A,i}.
\end{align}
One can verify that this closed real semi-positive $(1,1)$-form is independent of the choice of $A$ and the orthogonal coordinates $(x_1,\ldots,x_N)$ for $V$ (see \cite[\S 3]{DW24b}). Moreover, it is invariant under the $\pi_1(X)$-action. Consequently, it descends to a smooth, real, closed, semi-positive $(1,1)$-form on $\mathcal{R}(u)$. 
The Lipschitz property of $u$ and elliptic regularity imply  that this form extends to a positive closed $(1,1)$-current $T_{\varrho}$ on $X$ with continuous potential. 
\begin{dfn}[Canonical current]\label{def:canonical}
	The closed positive $(1,1)$-current $T_\varrho$ on $X$ defined above is called the \emph{canonical current} of $\varrho$.
\end{dfn}

\begin{thm}[{\cite[Theorem E]{CDY22}}]\label{thm:KE}
	Let $X$ be a  quasi-projective normal variety and let $\rho \colon \pi_1(X) \to \GL_n(K)$ be a   representation, where $K$ is a non-archimedean local field. Then there exists a  dominant morphism $s_\rho \colon X \to S_\rho$ to a normal projective variety with connected general fibers, such that for any irreducible closed subvariety $Z \subset X$, the following properties are equivalent: 
	\begin{thmlist}
		\item \label{item:KZ1} The image $\rho(\operatorname{Im}[\pi_1(Z^{\mathrm{norm}}) \to \pi_1(X)])$ is bounded;
		\item \label{item:KZ3} The image $\rho(\operatorname{Im}[\pi_1(Z) \to \pi_1(X)])$ is bounded;
		\item \label{item:KZ2} The image $s_\rho(Z)$ is a point;
		\item \label{item:KZ4} The restriction of the canonical current $T_\rho|_Z$ defined in \cref{def:canonical} is trivial.
		\item  \label{item:KZ5} The restriction of the multivalued 1-form $\eta_\rho|_Z$ defined in \cref{def:multivaluedform} is trivial.
	\end{thmlist}   
\end{thm} 

\begin{lem}\label{lem:period}
	If   $\varrho:\pi_1(X)\to \GL_{N}(\bC)$ underlies a $\bC$-VHS $\cL$ with discrete monodromy, then after replacing $X$ by a partial compactification, the $\bC$-VHS extends, and the associated period map $p:X\to \sD/\Gamma$ is proper, where $\sD$ is the period domain and $\Gamma$ is the monodromy group of $\cL$. In this case, the Shafarevich morphism ${\rm sh}_\varrho:X\to {\rm Sh}_\varrho(X)$ of $\varrho$ is the Stein factorization of $p$.
\end{lem}

\begin{proposition}[{\cite[Proposition 3.13]{DY23}}]\label{prop:cons}
	There exists a $\bC$-VHS $\cL$ on $X$ such that for any morphism $f:Y\to X$ from a smooth quasi-projective variety $Y$ to $X$, if $s_{\rm fac}\circ f(Y)$ is a point, then for any semisimple representation $\varrho:\pi_1(X)\to \GL_{N}(\bC)$,  $f^*\varrho$ is a direct factor of the $\bC$-VHS $f^*\cL$, which has discrete monodromy.
\end{proposition}

\begin{lem}[{\cite[Lemma 3.1]{DY23b}}]\label{lem:finite group}
	Let $K$ be an algebraically closed field of positive characteristic and let $\Gamma$ be a finitely generated group. Let $\varrho:\Gamma\to \mathrm{GL}_N(K)$ be a representation such that its semisimplification has finite image. Then $\varrho(\Gamma)$ is finite.
\end{lem}

\begin{thm}[{\cite[Theorem 3.9]{DY23b}}]\label{thm:Shapositive} 
	Let $X$ be a smooth quasi-projective variety. Assume that there exists an almost faithful representation $\rho \colon \pi_1(X) \to \operatorname{GL}_{N}(K)$, where $K$ is a field of   characteristic $p>0$.	The Shafarevich morphism ${\rm sh}_X:X\to {\rm Sh(X)}$ of $X$  is obtained through the simultaneous Stein factorization $s_\infty:X\to S_\infty(X)$  of the reductions   $\{s_{\tau}:X\to S_\tau\}_{\tau\in \Upsilon}$,  where
	$\Upsilon$ consists of all reductive representations with $K$  a non-archimedean local field of characteristic $p$,  and $s_\tau:X\to S_\tau$ is the Katzarkov-Eyssidieux reduction map  defined in \cref{thm:KE}.  
\end{thm} 

\begin{thm}[{\cite[Theorem B]{BDDM}}]\label{thm:sym}
	Let $X$ be a smooth quasi-projective variety. If there exists a repersentation $\varrho:\pi_1(X)\to\GL_{N}(K)$ for any field $K$ such that $\varrho(\pi_1(X))$ is infinite. Then  for some $k\in \bN$, 
	we have
	$$
	H^0(\overline{X}, {\rm Sym}^k\Omega_{\overline{X}}(\log D))\neq 0,
	$$
	where  $\overline{X}$ is a smooth projective compactification, such that $D=\overline{X}\backslash X$ is a simple normal crossing divisor. 
\end{thm}   

\begin{proof}[Proof  of \cref{thm:sym} (sketch)]
	The existence of logarithmic symmetric differentials is preserved under finite \'etale covers.
	Therefore, throughout the proof, we may freely replace $X$ by a finite \'etale cover.
	
	\medskip
	\noindent\textbf{Case~1:} ${\rm char}\, K = 0$.
	
	We may assume that $K=\bC$. Let $\sigma$ be the semisimplification of $\varrho$.
	
	\medskip
	\noindent\textbf{Case~1.1:} $\sigma$ has finite image.
	
	After replacing $X$ by a finite \'etale cover, the representation $\varrho$ becomes unipotent.
	Its image then admits a filtration whose successive quotients are abelian. Consequently,
	there exists an abelian representation of $\pi_1(X)$ with infinite image. This implies that
	$H^1(X,\bC)$ is infinite-dimensional.
	
	Recall that
	\[
	H^1(X,\bC)
	= H^0\bigl(\overline{X}, \Omega_{\overline{X}}(\log D)\bigr)
	\oplus H^{0,1}(X).
	\]
	Hence $H^0\bigl(\overline{X}, \Omega_{\overline{X}}(\log D)\bigr)\neq 0$.
	
	\medskip
	\noindent\textbf{Case~1.2:} $\sigma$ has infinite image.
	
	\medskip
	\noindent\textbf{Case~1.2.1:}  
	For every representation $\tau:\pi_1(X)\to \GL_N(K)$, where $K$ is a non-archimedean local
	field of characteristic zero, the image of $\tau$ is bounded.
	
In this case, the factorization map \(s_{\rm fac}: X \to S_{\rm Fac}(X)\) is constant. 
By \cref{prop:cons}, \(\sigma\) appears as a direct factor of a complex variation of Hodge 
structure \(\mathcal{L}\) with discrete monodromy. 
Since \(\sigma\) has infinite image, it follows that the monodromy of \(\mathcal{L}\) is infinite.
	
	By \cref{lem:period}, the associated period map $p:X\to \sD/\Gamma$ has positive-dimensional
	image. Let $Z$ be a desingularization of the closure of $p(X)$. Using curvature properties
	of period domains in the horizontal direction \cite{Gri70}, Brunebarbe and Cadorel
	\cite{BC20} proved that $Z$ has a big logarithmic cotangent bundle. In particular,
	$Z$ admits many logarithmic symmetric differentials. Pulling them back to $X$ yields
	the desired conclusion.
	
	\medskip
	\noindent\textbf{Case~1.2.2:}  
	There exists a representation $\tau:\pi_1(X)\to \GL_N(K)$, where $K$ is a non-archimedean
	local field of characteristic zero, whose image is unbounded.
	
By \cref{thm:KE}, there exists a nontrivial multivalued logarithmic \(1\)-form \(\eta_{\tau}\).
Taking local products of \(\eta_{\tau}\), one obtains logarithmic symmetric differentials on a 
Zariski open subset of \(\overline{X}\).
One then shows that these extend to logarithmic symmetric differentials on the pair 
\((\overline{X}, D)\).
	
	\medskip
	
	\noindent\textbf{Case~2:} ${\rm char}\, K = p>0$.
	
	By the proof of \cref{thm:Shapositive}, it follows that there exists an unbounded
	representation
	\[
	\tau:\pi_1(X)\to \GL_N\bigl(\bF_q((t))\bigr),
	\]
	where $q=p^\ell$ for some $\ell\in \bN$. Otherwise, the character variety
	$M_{\rm B}(X,N)_{\bF_p}$ would be zero-dimensional, and it would follow from
	\cref{lem:finite group} that $\varrho$ has finite image, contradicting the assumption.
	We then apply the same argument as in  {Case~1.2.2} to conclude that $X$
	admits a non-trivial logarithmic symmetric differential.
\end{proof}

\begin{thebibliography}{99}
\bibitem[BC20]{BC20} Yohan Brunebarbe and Beno{\^{\i}}t Cadorel, \emph{Hyperbolicity of varieties supporting a variation of {Hodge} structure}, Int. Math. Res. Not. \textbf{2020} (2020), no.~6, 1601--1609 (English).
\bibitem[BDDM]{BDDM} Damian {Brotbek}, Georgios {Daskalopoulos}, Ya~{Deng}, and Chikako {Mese}, \emph{{Pluriharmonic maps into buildings and symmetric differentials}}, arXiv e-prints (2022), arXiv:2206.11835.
\bibitem[CDY22]{CDY22} \bysame, \emph{{Hyperbolicity and fundamental groups of complex quasi-projective varieties (II): via non-abelian Hodge theories}}, arXiv e-prints (2025), arXiv:2512.15367.
\bibitem[DM24]{DM24} Ya~{Deng} and Chikako {Mese}, \emph{Existence and unicity of pluriharmonic maps to euclidean buildings and applications}, New Aspects of Teichmüller Theory, Adv. Stud. Pure Math., Math. Soc. Japan, Tokyo, 2026, to appear.
\bibitem[DW24b]{DW24b} Ya~{Deng} and Botong {Wang}, \emph{{$L^2$-vanishing theorem and a conjecture of Koll{\'a}r}}, arXiv e-prints (2024), arXiv:2409.11399.
\bibitem[DY23]{DY23} Ya~{Deng}, Katsutoshi {Yamanoi}, and Ludmil {Katzarkov}, \emph{{Reductive Shafarevich Conjecture}}, arXiv e-prints (2023), arXiv:2306.03070.
\bibitem[DY23b]{DY23b} Ya~{Deng} and Katsutoshi {Yamanoi}, \emph{Linear shafarevich conjecture in positive characteristic, hyperbolicity and applications}, J. Reine Angew. Math. (2025).
\bibitem[Gri70]{Gri70} Phillip~A. Griffiths, \emph{Periods of integrals on algebraic manifolds. {III}. {S}ome global differential-geometric properties of the period mapping}, Inst. Hautes \'{E}tudes Sci. Publ. Math. (1970), no.~38, 125--180. \MR{282990}
\end{thebibliography}
\end{document}
